\chapter{系统总体设计}

\section{系统架构设计}

\subsection{前后端分离架构概述}

在整体架构的选型上，本系统选择了前后端分离这一技术路线。具体而言，前端部分以 Vue~3 来构建单页面应用（SPA），由用户浏览器负责运行；后端则基于 Node.js + Express 来搭建 RESTful API 服务层；二者之间以 HTTP 协议开展数据交互，而 AI 流式对话这一特殊通道，则选用了 Server-Sent Events（SSE）机制来加以实现。选择前后端分离的根本原因在于，这种方式能够把界面渲染逻辑与业务数据处理彻底区分开来，前后端各自可以独立开发、独立部署，后续如果只需要对某一端进行调整，也不必牵动整个项目的其他部分。

如果从纵向分层的视角来看，整个系统可以被划分成三个层次：最上层是\textbf{表示层}（前端 Vue 应用），主要承担页面渲染、用户交互以及前端状态管理；中间层是\textbf{业务逻辑层}（Express API 服务），负责对请求进行鉴权、执行业务规则并与数据库完成读写交互；最底层则是\textbf{数据持久层}（MySQL 数据库），主要用于对结构化数据进行长期存储。

三类用户——学生、咨询师以及管理员——共用同一套后端服务，借助不同的路由前缀和角色守卫中间件来实现彼此之间的功能隔离。这种\textbf{``单一后端、多端服务''}的设计思路，不仅避免了重复建设不同后端的麻烦，还从根本上保证了各端数据来自统一数据源，减少了由数据不同步引发问题的可能性。

\begin{figure}[H]
    \centering
    \includegraphics[width=0.85\textwidth]{fig3-1.png}
    \caption{系统整体架构图}
    \label{fig:3-1}
\end{figure}

\subsection{技术栈选型说明}

\textbf{前端：Vue~3 + Composition API。}相较于 Options API 依赖"选项合并"的逻辑复用方式，Vue~3 的 Composition API 将其转变为基于函数的组合模式，使复杂组件内部的逻辑能够围绕关注点拆分为独立且可复用的函数（Composable），而不必再被分散到 \texttt{data}、\texttt{computed}、\texttt{methods} 等选项之中\cite{vuejs3docs}。对于本系统中预约状态轮询、AI 流式消息处理等逻辑相对复杂的场景，Composition API 能够使代码组织更加清晰，并在一定程度上降低测试与维护成本。

\textbf{状态管理：Pinia。}Pinia 作为 Vue~3 官方推荐的状态管理库\cite{piniadocs}，相较于 Vuex~4，其 API 更为简洁，并且原生支持 TypeScript 类型推导，同时允许 Store 之间直接相互引用。系统通过 Pinia 对认证状态中的当前用户信息以及 JWT Token 进行统一管理，并承担跨组件共享业务数据的功能。

\textbf{后端：Node.js + Express。}Node.js 凭借事件驱动机制以及非阻塞 I/O 模型，在面对大量并发 I/O 请求时仍可以保持较高的吞吐量\cite{cantelon2018}，所以会更适宜 SSE 长连接这一类场景——服务端能够在不阻塞主线程的情况下，把 DeepSeek API 的流式响应实时转发到多个前端连接\cite{hahn2016}。Express 框架本身相对轻量，同时中间件生态也比较成熟，因此可以用较低的配置成本来搭建结构清晰的 RESTful 服务。

\textbf{ORM：Sequelize~6。}Sequelize 会把数据库操作封装成 JavaScript 对象以及方法，还支持 async/await 语法，这样一来，业务代码就可以较好地摆脱原始 SQL 所带来的繁琐性和安全隐患\cite{seqdocs}。借助在模型层统一设置 \texttt{underscored: true}，JavaScript 端的 camelCase 字段名以及数据库端的 snake\_case 列名之间能够实现自动的双向映射，进而让两端的命名风格分别保持一致。

\textbf{认证：JWT。}JSON Web Token 具有无状态特性，服务端无需维护 Session 存储，因此更适合进行水平扩展；同时，Token 的 Payload 中包含用户 id 与 role 字段，使角色守卫中间件无需额外查询数据库即可完成权限判定，从而降低认证链路上的数据库压力\cite{jones2015}。


\section{数据库设计}

\subsection{总体设计原则}

数据库方面选用 MySQL~8，字符集统一进行配置，设为支持 Emoji 以及完整 CJK 字符集的 utf8mb4，排序规则则设置成 \texttt{utf8mb4\_unicode\_ci}\cite{mysql80docs}。同时，所有数据表都会保留 \texttt{created\_at} 以及 \texttt{updated\_at} 两个时间戳字段，并且由 Sequelize 自动维护，用来记录数据的创建时间以及最后修改时间。

\textbf{系统共设计 10 张数据表，按职责可分为以下三组：}
\begin{itemize}
    \item \textbf{用户与档案组：}\texttt{users}、\texttt{student\_profiles}、\texttt{counselor\_profiles}
    \item \textbf{咨询业务组：}\texttt{appointments}、\texttt{sessions}、\texttt{session\_notes}、\texttt{feedback}
    \item \textbf{辅助功能组：}\texttt{announcements}、\texttt{ai\_conversations}、\texttt{ai\_messages}
\end{itemize}

图~\ref{fig:3-ea}~给出系统各主要实体的属性图，清晰展示每张核心数据表所包含的属性字段。

% ---- 实体属性图（6 个主要实体） ----
\begin{figure}[H]
\centering
\resizebox{\linewidth}{!}{%
\begin{tikzpicture}[
    entity/.style={draw, fill=blue!15, rectangle,
                   minimum width=2.6cm, minimum height=0.65cm,
                   text centered, font=\heiti\small},
    attr/.style={draw, fill=gray!8, ellipse,
                 minimum width=2.0cm, minimum height=0.5cm,
                 text centered, font=\scriptsize, inner sep=1pt},
    line/.style={thin, gray}
]
%---------- users ----------
\node[entity] (U) at (0,0) {用户表 (users)};
\node[attr] (Ua) at (-2.8, 0.9)  {id (PK)};
\node[attr] (Ub) at (-1.0, 1.5)  {username};
\node[attr] (Uc) at ( 1.0, 1.5)  {email};
\node[attr] (Ud) at ( 2.8, 0.9)  {password\_hash};
\node[attr] (Ue) at ( 2.4,-0.9)  {role};
\node[attr] (Uf) at (-2.4,-0.9)  {is\_active};
\foreach \x in {Ua,Ub,Uc,Ud,Ue,Uf} \draw[line] (U)--(\x);

%---------- appointments ----------
\node[entity] (A) at (8.5,0) {预约表 (appointments)};
\node[attr] (Aa) at (6.0, 1.0)  {id (PK)};
\node[attr] (Ab) at (7.5, 1.6)  {student\_id};
\node[attr] (Ac) at (9.5, 1.6)  {counselor\_id};
\node[attr] (Ad) at (11.0, 1.0) {appointment\_date};
\node[attr] (Ae) at (11.2,-0.8) {status};
\node[attr] (Af) at (9.2,-1.4)  {reason};
\node[attr] (Ag) at (7.2,-1.4)  {cancel\_reason};
\node[attr] (Ah) at (5.8,-0.8)  {cancelled\_by};
\foreach \x in {Aa,Ab,Ac,Ad,Ae,Af,Ag,Ah} \draw[line] (A)--(\x);
\end{tikzpicture}%
}% end resizebox

\vspace{0.6em}

\resizebox{\linewidth}{!}{%
\begin{tikzpicture}[
    entity/.style={draw, fill=blue!15, rectangle,
                   minimum width=2.6cm, minimum height=0.65cm,
                   text centered, font=\heiti\small},
    attr/.style={draw, fill=gray!8, ellipse,
                 minimum width=2.0cm, minimum height=0.5cm,
                 text centered, font=\scriptsize, inner sep=1pt},
    line/.style={thin, gray}
]
%---------- sessions ----------
\node[entity] (S) at (0,0) {会话表 (sessions)};
\node[attr] (Sa) at (-2.5, 1.0) {id (PK)};
\node[attr] (Sb) at (-0.5, 1.6) {appointment\_id};
\node[attr] (Sc) at ( 1.8, 1.6) {counselor\_id};
\node[attr] (Sd) at ( 2.8, 0.2) {started\_at};
\node[attr] (Se) at ( 2.2,-1.0) {ended\_at};
\node[attr] (Sf) at (-2.2,-1.0) {status};
\foreach \x in {Sa,Sb,Sc,Sd,Se,Sf} \draw[line] (S)--(\x);

%---------- counselor_profiles ----------
\node[entity] (CP) at (8.5,0) {咨询师档案 (counselor\_profiles)};
\node[attr] (CPa) at (6.2, 1.0) {id (PK)};
\node[attr] (CPb) at (7.8, 1.6) {user\_id};
\node[attr] (CPc) at (9.8, 1.6) {title};
\node[attr] (CPd) at (11.0, 0.4) {specialties};
\node[attr] (CPe) at (10.8,-0.9) {available\_slots};
\node[attr] (CPf) at (8.0,-1.4) {bio};
\node[attr] (CPg) at (6.2,-0.9) {real\_name};
\foreach \x in {CPa,CPb,CPc,CPd,CPe,CPf,CPg} \draw[line] (CP)--(\x);
\end{tikzpicture}%
}% end resizebox

\vspace{0.6em}

\resizebox{\linewidth}{!}{%
\begin{tikzpicture}[
    entity/.style={draw, fill=blue!15, rectangle,
                   minimum width=2.6cm, minimum height=0.65cm,
                   text centered, font=\heiti\small},
    attr/.style={draw, fill=gray!8, ellipse,
                 minimum width=2.0cm, minimum height=0.5cm,
                 text centered, font=\scriptsize, inner sep=1pt},
    line/.style={thin, gray}
]
%---------- ai_conversations ----------
\node[entity] (AC) at (0,0) {AI 对话 (ai\_conversations)};
\node[attr] (ACa) at (-2.6, 1.0) {id (PK)};
\node[attr] (ACb) at (-0.4, 1.6) {student\_id};
\node[attr] (ACc) at ( 2.0, 1.6) {title};
\node[attr] (ACd) at ( 2.8,-0.2) {updated\_at};
\node[attr] (ACe) at (-2.8,-0.4) {created\_at};
\foreach \x in {ACa,ACb,ACc,ACd,ACe} \draw[line] (AC)--(\x);

%---------- ai_messages ----------
\node[entity] (AM) at (8.5,0) {AI 消息 (ai\_messages)};
\node[attr] (AMa) at (6.4, 1.0) {id (PK)};
\node[attr] (AMb) at (8.0, 1.6) {conversation\_id};
\node[attr] (AMc) at (10.0, 1.6) {role};
\node[attr] (AMd) at (11.0, 0.0) {content};
\node[attr] (AMe) at (6.4,-0.8) {created\_at};
\foreach \x in {AMa,AMb,AMc,AMd,AMe} \draw[line] (AM)--(\x);
\end{tikzpicture}%
}% end resizebox
    \caption{系统主要实体属性图}
    \label{fig:3-ea}
\end{figure}

图~\ref{fig:3-2}~展示了 10 张数据表之间的完整关联关系（ER 图）。

\begin{figure}[H]
    \centering
    \includegraphics[width=\textwidth]{fig3-2.png}
    \caption{数据库 ER 图}
    \label{fig:3-2}
\end{figure}

\subsection{核心表结构说明}

\textbf{users 表。}系统中的用户中心表，会由全部角色共同来使用。它的核心字段包括：唯一索引字段 \texttt{username}（长度限制在 2—50 个字符之间）、唯一索引字段 \texttt{email}、\texttt{password\_hash}（bcrypt 哈希值）、\texttt{role}（ENUM：\texttt{student} / \texttt{counselor} / \texttt{school}），以及 TINYINT 字段 \texttt{is\_active}，用来控制账号是否处于可用状态。\texttt{role} 字段是权限控制的核心依据，在 JWT 生成时把它写入 Payload 之后，中间件在验证阶段就可以直接读取，所以不需要再回查数据库。

\textbf{student\_profiles 与 counselor\_profiles 表。}这两张档案扩展表主要用来承载用户档案方面的扩展信息，都会借助 \texttt{user\_id} 与 \texttt{users} 表建立一对一的关联关系。\texttt{counselor\_profiles} 中的 \texttt{specialties} 和 \texttt{available\_slots} 字段都选用 JSON 类型来进行存储，前者用来保存擅长方向的标签数组，后者用来保存可预约时段的配置内容，这样就能够较为灵活地适配不同咨询师的排班需求，也不必再为这类可变结构额外设计关联表。

\textbf{appointments 表。}预约核心表，主要用来记录每一条预约从创建到结束的完整生命周期。\texttt{appointment\_date} 字段选用 \texttt{DATEONLY} 类型（对应 MySQL \texttt{DATE}），同时在 Sequelize 配置里启用 \texttt{dateStrings: ['DATE']}，让驱动把日期列直接按 \texttt{'YYYY-MM-DD'} 的字符串形式返回，这样就能够避免因时区换算而产生日期偏移的问题。\texttt{status} 字段的枚举值设为 \texttt{pending}、\texttt{confirmed}、\texttt{cancelled}、\texttt{completed} 以及 \texttt{no\_show}，用来覆盖预约流程中的完整状态流转。

\textbf{sessions 表。}咨询会话记录表，与 \texttt{appointments} 表形成一对一的关联关系（\texttt{appointment\_id} 设有唯一索引）。\texttt{started\_at} 以及 \texttt{ended\_at} 分别用来记录实际会话开始和结束的时间，两者之间形成的时间差，也就是本次咨询的时长，可以为后续开展统计分析工作提供依据。

\textbf{session\_notes 表。}案例笔记表，与 \texttt{sessions} 表形成一对多的关联关系，这样就允许咨询师围绕同一会话持续补充多条笔记。借助 \texttt{counselor\_id} 外键来对数据归属进行明确限制，并且在服务端查询时加上 \texttt{counselor\_id = req.user.id} 这一条件，就可以在 SQL 层面防止出现跨咨询师的数据越权访问。

\textbf{ai\_conversations 与 ai\_messages 表。}AI 对话的两级数据结构。其中，\texttt{ai\_conversations} 表对应一次完整的会话单元，\texttt{ai\_messages} 表负责存放该会话中的逐条消息内容，\texttt{role} 字段则用来区分用户输入（\texttt{user}）以及模型回复（\texttt{assistant}）。在构建历史上下文时，服务端会查询最近 20 条 \texttt{ai\_messages} 记录，并按照时间顺序把它们组装成消息数组后发送给 DeepSeek API，这样一来，模型就能够感知对话的连续性。

\subsection{外键约束与索引设计}

系统在模型层借助 Sequelize 所提供的关联方法（\texttt{hasOne}、\texttt{hasMany}、\texttt{belongsTo}）来声明外键关系，ORM 在同步数据库结构时会自动生成对应的外键约束，进而保证引用完整性。

{\sloppy
高频查询字段的索引设计可以概括为：\texttt{users.username} 以及 \texttt{users.email} 会由于 \texttt{unique:~true} 这一约束而自动生成唯一索引；\texttt{appointments.student\_id}、\texttt{appointments.counselor\_id} 以及 \texttt{appointments.status} 设置普通索引，用来支撑预约列表当中的条件过滤查询；\texttt{ai\_messages.conversation\_id} 也建立索引，从而加快对话历史的分页查询\cite{schwartz2013}。
}

综合来看，外键约束以及索引策略共同在数据库层面形成了兼顾完整性与性能的双重保障。前者会在数据写入阶段对引用关系开展约束工作，从而防止出现孤立记录，并且保证引用完整性；后者则会在数据读取阶段对高频查询进行加速，进而降低发生全表扫描的概率。二者相互配合之后，使系统即便在数据量持续增长的情况下，仍然可以保持较为稳定的查询响应时间，并为上层业务逻辑的正确运行提供可靠的数据层支持。


\section{系统功能模块划分}

系统按功能职责划分为若干顶层模块，各模块内部再细化为具体的子功能单元，从而形成层次较为清晰的整体结构。

\begin{figure}[H]
    \centering
    \includegraphics[width=0.90\textwidth, trim=0 2.0cm 0 2.0cm, clip]{fig3-3.png}
    \caption{系统功能模块层次图}
    \label{fig:3-3}
\end{figure}

如图~\ref{fig:3-3}~所示，系统按照功能职责来进行划分，被分为认证、学生端、咨询师端、学校管理端、AI 心理助手以及公共服务六个顶层模块。其中，认证模块作为安全基础，会贯穿在各端的路由当中；各业务端模块围绕对应用户角色的核心使用场景来展开；AI 心理助手模块则借助 SSE 流式对话方式，独立提供智能心理支持；公共服务模块负责对错误处理、日志记录等可跨端复用的基础能力进行封装，从而保证系统在运行过程中具备统一保障。各顶层模块内部还会进一步细分为具体的子功能单元，模块之间依靠统一的 RESTful API 接口以及中间件机制来实现低耦合调用与协作。以下各小节将逐一对各模块的核心职责及功能边界进行阐述。

\subsection{认证模块}

认证模块作为整个系统安全体系的基础部分，负责对用户身份开展建立与验证工作。注册接口会接收用户名、邮箱、密码以及角色信息，并在把密码进行 bcrypt 哈希处理后写入数据库；登录接口在完成凭据校验之后会生成包含 \texttt{id} 和 \texttt{role} 字段的 JWT，它的有效期被设置为 7 天；\texttt{authMiddleware} 中间件会在每一次受保护请求中提取并验证 Token，\texttt{roleGuard} 还会在此基础上进一步开展角色权限校验。该模块以中间件方式无侵入地嵌入到各端路由当中，使业务逻辑无需额外去处理鉴权细节。

\subsection{学生端模块}

学生端模块主要包括个人档案管理、咨询师浏览、预约管理、公告查阅以及咨询评价五个子模块。各子模块都会借助 \texttt{authMiddleware} + \texttt{roleGuard('student')} 来形成双重保护机制，从而确保只有携带学生身份 Token 的请求才可以访问。与此同时，预约管理子模块还会和咨询业务模块中的 \texttt{appointments} 数据层保持联动，使学生发起的预约申请能够实时进入咨询师端的待审核队列。

\subsection{咨询师端模块}

咨询师端模块主要围绕预约审核以及会话管理两项核心工作来展开。预约审核子模块负责对来访申请进行确认或者拒绝；会话管理子模块会在预约被确认之后由咨询师手动创建会话记录，并且在会谈结束时执行"结束会话"操作，再由系统自动写入结束时间；案例笔记子模块和会话之间形成一对多关联，因此允许咨询师在会话进行期间或者结束之后随时补充记录。

\subsection{学校管理端模块}

管理端模块主要承担平台运营层面的相关职责，因此通常不会直接介入具体的咨询内容。咨询师账号管理子模块具有创建账号，以及对账号进行启用、禁用的功能；数据统计子模块会对平台注册用户数、预约量、完成率、AI 对话次数以及平均评分等关键指标开展汇总，并以结构化数据的形式返回，以便前端去完成图表渲染；公告管理子模块支持对公告内容进行发布与删除，相关公告会在学生端以及咨询师端同步可见。

\subsection{AI 心理助手模块}

{\sloppy
AI 助手模块在系统架构层面上独立于三端的业务路由，并且借助专属的 \texttt{/api/v1/student/ai} 路由来对外提供服务。模块内部把有关于 DeepSeek API 的全部交互流程统一封装成一条完整链路：接收用户消息 $\rightarrow$ 从 \texttt{ai\_messages} 表当中加载最近 20 条上下文 $\rightarrow$ 完成请求体的拼装 $\rightarrow$ 以流式方式调用 DeepSeek API $\rightarrow$ 把响应块实时转发到前端 SSE 连接 $\rightarrow$ 将完整回复写入数据库。该模块依赖 \texttt{services/ai.service.js} 来对 DeepSeek API 调用进行封装，其他模块不得绕过 Service 层直接访问外部 AI 接口，这样可以保证调用入口的唯一性，以及后续维护工作的便利性\cite{deepseek2024}。
}


\section*{本章小结}

本章从系统架构、数据库以及功能模块这三个层面，对系统的总体设计进行了说明。在架构层面，确立了以 Vue~3 + Node.js/Express + MySQL 为核心的前后端分离技术路线，同时对各项技术选型的合理性作了说明；在数据库层面，完成了 10 张数据表的结构设计，明确了核心表字段、外键约束以及索引策略，还完整呈现了 ER 关系；在功能模块层面，把系统划分为认证、学生端、咨询师端、管理端以及 AI 助手五大模块，并且阐明了各模块的职责边界以及内部结构。本章所形成的设计结果将直接为第 5 章的详细实现提供依据。
